\documentclass[runningheads]{llncs}
%
\usepackage{tikz}
\usetikzlibrary{positioning, arrows.meta}
% cite package removed: llncs handles citations natively with splncs04.bst
\usepackage{booktabs}
\usepackage{multirow}
\usepackage{makecell}
\usepackage{graphicx}
\usepackage{amsmath,amssymb,amsfonts}
\usepackage{algorithmic}
\usepackage{textcomp}
\usepackage{xcolor}
\usepackage{enumitem}
\setlist{noitemsep,topsep=2pt}

\begin{document}

\title{Siamese-based Real-ESRGAN with High-Order Degradation for Realistic Image Restoration}
\titlerunning{Siamese-based Real-ESRGAN for Realistic Image Restoration}

\author{Le Hoai Nam\inst{1} \and Doan Nhat Quang\inst{1} \and Giang Anh Tuan\inst{1} \and\\
Tran Hoang Tung\inst{1} \and Le Huu Ton\inst{2} \and Nguyen Hoang Ha\inst{1}}

\authorrunning{L. H. Nam et al.}

\institute{University of Science and Technology of Hanoi, Vietnam Academy of Science and Technology, Hanoi, Vietnam\\
\email{nguyen-hoang.ha@usth.edu.vn} \and
CMC University, Hanoi, Vietnam}

\maketitle

\begin{abstract}
The application of deep learning to enhance image quality is a rapidly emerging research area. However, restoring realistic images with complex degradations such as noise, blur, and JPEG compression remains a major challenge. To address this, we propose a novel Siamese-based framework that enhances the Real-ESRGAN architecture through a combination of realistic degradation modeling and dual-branch knowledge distillation. First, we implement a High-order Realistic Degradation pipeline in the data synthesis stage, utilizing randomized operations and explicit camera noise simulation to bridge the sim-to-real gap. Furthermore, we introduce a Siamese training strategy consisting of a Teacher network and a Student network with same layers. In this framework, the two branches process different degradation levels of the same ground-truth image: the Teacher learns from ``lightly'' degraded samples to provide high-quality structural guidance, while the Student is trained on ``severely'' degraded inputs. By distilling knowledge from the Teacher branch to the Student, our model effectively restores fine details while maintaining low computational costs. Experimental results show that our Siamese approach achieves superior perceptual quality, specifically a NIQE score of 3.82, outperforming the baseline in stability and visual realism. Our framework offers a promising solution for real-world tasks such as enhancing old photographs and surveillance imagery.

\keywords{Image Super-Resolution \and Siamese Network \and Knowledge Distillation \and Realistic Degradation \and Real-ESRGAN.}
\end{abstract}


\section{Introduction}

Image restoration aims to reconstruct high-quality images from complex real-world degradations. While diffusion models and Transformers achieve strong performance, their high computational costs limit real-time deployment. GAN-based methods like ESRGAN~\cite{wang2018esrgan} and BSRGAN~\cite{zhang2021designing} offer efficient alternatives. Real-ESRGAN~\cite{wang2021realesrgan} further improved blind super-resolution via high-order degradation modeling. However, most existing methods rely on a single degradation path, limiting robustness.

To address this issue, we propose a Siamese-based Real-ESRGAN framework with High-Order Degradation for realistic image restoration. This work builds upon and significantly extends our previous research in~\cite{11309534}. While our prior work focused on basic knowledge distillation for model compression, this paper introduces a novel weight-sharing Siamese architecture and a curriculum learning strategy to specifically tackle the sim-to-real gap in heavy degradation scenarios.

Our main contributions are summarized as follows:
\begin{itemize}
    \item We propose a Weight-Sharing Siamese (Twin) architecture that enables implicit co-evolution of features, allowing the model to learn robust representations from multi-level degradation pairs.
    \item We implement a three-phase Curriculum Learning strategy that progressively increases degradation difficulty, stabilizing the training process for complex real-world artifacts.
    \item We refine the Realistic High-Order Degradation pipeline with randomized operation ordering and explicit camera-noise simulation, significantly improving data diversity.
\end{itemize}

\section{Related Work}

Early work in image super-resolution (SR) began with SRCNN~\cite{dong2014learning}, which established the foundation of deep learning-based image restoration. To improve perceptual quality, GAN-based methods such as ESRGAN~\cite{wang2018esrgan} were later introduced, while residual-based architectures like EDSR~\cite{lim2017enhanced} and RDN~\cite{zhang2018residual} demonstrate the effectiveness of deep residual learning. However, most early methods rely on bicubic downsampling to generate low-quality (LQ) images, which cannot accurately represent real-world degradations.

Research then shifted toward Blind Super-Resolution (Blind SR)~\cite{zhang2018learning,luo2020unfolding,gu2019blind}, which aims to handle unknown degradation patterns such as blur, compression, and sensor noise. Real-ESRGAN~\cite{wang2021realesrgan} addressed this problem by introducing a high-order degradation model that generates more realistic synthetic training pairs, significantly improving robustness on real-world images. In our previous work~\cite{11309534}, we further improved this degradation process using randomized operation order, refined parameter ranges, and explicit camera-noise simulation.

To improve deployment efficiency, Knowledge Distillation (KD)~\cite{hinton2015distilling,Hui_2019} has been widely adopted to compress large restoration models into lightweight networks. Recent methods such as FAKD~\cite{gao2022featuredistillationinteractionweighting} utilize feature-aware distillation to preserve structural details. Conventional KD mainly relies on output-level supervision from a fixed Teacher. In contrast, our work introduces a Siamese Knowledge Distillation framework, where Teacher and Student (sharing the same weights) jointly learn from different degradation levels, enabling stronger feature consistency.


\noindent\textbf{Siamese Networks.} Originally proposed for signature verification~\cite{bromley1993signature}, Siamese networks have been widely applied to similarity learning tasks~\cite{koch2015siamese}. We adapt this concept to image restoration by processing pairs of images with different degradation severities through a single shared-weight generator.

\noindent\textbf{Curriculum Learning.} Curriculum learning~\cite{bengio2009curriculum} mimics human learning by starting with easy examples and progressively increasing difficulty, improving generalization and convergence~\cite{soviany2022curriculum}. We integrate this strategy by progressively increasing degradation strength across three training phases.

\noindent\textbf{Diffusion-based Super-Resolution.} Recent models such as StableSR~\cite{wang2024exploiting}, DiffBIR~\cite{lin2023diffbir}, and PASD~\cite{yang2023pixel} achieve state-of-the-art perceptual quality by leveraging generative priors, but suffer from high computational costs and slow inference. Our GAN-based Siamese approach targets comparable realism at significantly lower overhead.


\section{Methodology}

This study proposes an image restoration framework that combines an improved Realistic High-Order Degradation process with a Siamese joint training architecture and curriculum learning strategy. Instead of conventional Teacher--Student distillation, our method jointly trains dual branches using different degradation levels of the same image. The objective is to improve robustness under complex real-world degradations while maintaining restoration quality and computational efficiency. An overview of the proposed pipeline is shown in Fig.~\ref{fig:workflow}.
\begin{figure*}[!t]
    \centering
    \includegraphics[width=\textwidth,height=3.2cm]{images/pipeline.png}
    \caption{Overview of the proposed Siamese-based Real-ESRGAN framework~\cite{11309534}.}
    \label{fig:workflow}
\end{figure*}

\subsection{Data Augmentation by Realistic High-Order Degradation Method}

Most existing image restoration methods, including Real-ESRGAN~\cite{wang2021realesrgan}, generate synthetic low-quality (LQ) images from high-quality (HQ) images using a two-stage degradation pipeline. Each stage applies resizing, blurring, noise injection, and compression in a fixed order. Although effective, this strategy still cannot fully represent the diversity of real-world degradations.

Our method in Fig.~\ref{fig:degradation} keeps the general two-stage process (resizing, blurring, noise addition, JPEG compression, and camera-specific noise) while introducing several improvements:

\begin{figure*}[!t]
    \centering
    \includegraphics[width=\textwidth,height=2.6cm]{images/Realdegradationstructure.png}
    \caption{Overview of the proposed two-stage realistic degradation pipeline.}
    \label{fig:degradation}
\end{figure*}

\begin{itemize}
    \item \textbf{Randomized Degradation Order:} Flexible degradation sequence instead of a fixed order, while keeping resizing first.
    \item \textbf{Refined Parameter Ranges:} JPEG compression limited to narrower ranges to avoid unrealistic artifacts.
    \item \textbf{Explicit Camera Noise:} Chromatic aberration and sensor-specific pattern noise simulate smartphone and low-quality camera artifacts.
    \item \textbf{Pre-Generated Degradation:} LQ--HQ pairs generated offline, improving reproducibility and training efficiency.
    \item \textbf{Multi-Level Degradation:} Four quality levels (90\%, 80\%, 70\%, 60\%) are generated per GT image to support the curriculum Siamese framework.
\end{itemize}

For $\times4$ super-resolution training, LQ images are generated at one-quarter the resolution of their GT counterparts. GT images are cropped into $480 \times 480$ patches with a stride of 240 (50\% overlap), while LQ images are cropped into $120 \times 120$ patches with a stride of 60, maintaining the $\times4$ scale ratio.

Unlike standard SR training with a single degradation level, our Siamese framework uses paired inputs with different degradation severities. In each phase, the Teacher branch receives a mildly degraded image, while the Student branch receives a more severely degraded one. The GT image serves as the reference for pixel-wise, perceptual, and adversarial losses.

\subsection{Weight-Sharing Siamese Architecture with Curriculum Learning}

After generating the training data using the realistic high-order degradation pipeline, we improve restoration robustness through a Siamese joint training strategy with curriculum learning. Unlike conventional Teacher--Student distillation, which requires a fixed pre-trained Teacher and a separate lightweight Student, our framework uses a shared-weight dual-branch design to process different degradation levels of the same image. This reduces memory overhead while encouraging stronger feature consistency across degradation severities.

The proposed model is built on the Real-ESRGAN RRDBNet backbone~\cite{wang2018esrgan,wang2021realesrgan}, consisting of 23 RRDB blocks with 64 feature channels. Instead of compressing the Teacher into a smaller Student as in our previous work~\cite{11309534}, both branches share the same network parameters and are executed in two forward passes: a Teacher pass and a Student pass. A U-Net discriminator with spectral normalization is used for adversarial supervision.

During each training step, two degraded inputs are sampled from the same GT image with different degradation severities. The Teacher branch receives the relatively cleaner image, while the Student branch receives the more severely degraded one. For example, the first phase uses 90\% and 80\% quality pairs, followed by 80\%/70\% and 70\%/60\% in later curriculum stages. This progressive schedule allows the model to learn from easier restoration tasks before adapting to harder degradations.

The overall Siamese training framework is illustrated in Fig.~\ref{fig:siamese_architecture}. The optimization is performed through two forward passes. First, the Teacher pass processes the mildly degraded input under \texttt{torch.no\_grad()}, producing the pseudo-target output ($THQ$) and intermediate features ($T_{feat}$). Second, the Student pass restores the severely degraded input and produces ($SHQ$) and ($S_{feat}$). This curriculum-based dual-branch design enables progressive learning across degradation levels and improves feature consistency for real-world image restoration.

\begin{figure*}[!t]
    \centering
    \includegraphics[width=\textwidth,height=3.2cm]{images/student_teacher_siamese.drawio.png}
    \caption{Overview of the proposed weight-sharing Siamese training framework with curriculum learning.}
    \label{fig:siamese_architecture}
\end{figure*}

To optimize the network, we combine reconstruction loss, perceptual supervision~\cite{johnson2016perceptual}, adversarial learning, and feature distillation into the following unified objective:

\[
\begin{aligned}
L_{total} =\;&
w_{pixel} \, L_{pixel} \\
&+ w_{perceptual} \, L_{perceptual} \\
&+ w_{GAN} \, L_{GAN} \\
&+ w_{distill}^{out} \, L_{distill}^{out} \\
&+ w_{distill}^{feat} \, L_{distill}^{feat}
\end{aligned}
\]

where $L_{pixel}$ denotes the pixel-wise reconstruction loss between the Student output and the GT image, $L_{distill}^{out}$ measures the output-level consistency between Student and Teacher predictions, and $L_{distill}^{feat}$ enforces feature alignment between intermediate representations of both branches. The perceptual loss is computed based on the VGG19 architecture~\cite{simonyan2014very}.

Only the Student branch receives gradient updates, while the Teacher outputs are detached and used as dynamic supervision. Since both branches share parameters, optimizing difficult restoration cases also implicitly improves performance on milder degradations. This enables better convergence and stronger perceptual quality compared to static distillation frameworks.


\section{Experiments}
\subsection{Data}
We use the same data sources as our previous work~\cite{11309534}, including DF2K, Unsplash, and RealSR. DF2K consists of DIV2K, Flickr2K, and OutdoorScene, providing diverse high-quality (HQ) images with rich textures and large resolution variations. We further include additional images from Unsplash and RealSR to improve scene diversity and real-world degradation coverage.

Unlike conventional super-resolution training that relies on bicubic downsampling, we retain only the HQ images and apply our realistic high-order degradation pipeline to generate synthetic low-quality (LQ) counterparts. For the Siamese framework, each GT image is further degraded into multiple quality levels, enabling curriculum learning across different restoration difficulties.

Model evaluation is conducted on Set5, Set14, and RealSR benchmarks using PSNR, SSIM, and NIQE for quantitative comparison.

\subsection{Implementation Details}
Our implementation follows Real-ESRGAN~\cite{wang2021realesrgan} with Siamese modifications. We use Adam ($\beta_1=0.9$, $\beta_2=0.999$) with initial learning rate $1\times10^{-4}$, decayed by MultiStepLR ($\gamma=0.5$) at 200k, 400k, 500k, and 700k iterations over 800k total steps with 10k warmup. Losses: pixel L1 ($w_{pixel}=2.0$), perceptual via VGG19~\cite{simonyan2014very,johnson2016perceptual} ($w_{perceptual}=1.0$), adversarial ($w_{GAN}=0.5$), distillation ($\lambda_{kd}^{out}=\lambda_{kd}^{feat}=0.15$) and consistency (0.075), activated after warmup. All experiments run on an NVIDIA RTX 3090. Since both branches share the same RRDBNet backbone ($\approx$16.7M parameters), inference cost matches the Real-ESRGAN baseline ($\approx$0.08\,s for a $720{\times}480$ image), far below diffusion-based alternatives.

\subsection{Three-Phase Curriculum Training Schedule}
A central novelty of our framework is the \textbf{progressive multi-level degradation curriculum}. Rather than training on a single fixed degradation strength, the Siamese pair is exposed to pairs of increasing difficulty across three phases (Table~\ref{tab:curriculum}). The quality index (\%) denotes composite degradation intensity, where 100\% is the original GT and lower values indicate stronger blur, noise, and compression. In each phase, the Teacher receives the milder input and provides dynamic pseudo-supervision, while the Student is challenged with the harder degradation. This progressive schedule prevents the training instability observed when directly optimizing on the most severe pairs.

\begin{table}[h]
\centering
\caption{Three-phase curriculum training schedule (Deg.\ quality: 100\% = GT; lower = harder).}
\label{tab:curriculum}
\small
\begin{tabular}{@{}ccc p{3.6cm}@{}}
\toprule
\textbf{Phase} & \textbf{Teacher} & \textbf{Student} & \textbf{Training Focus} \\
\midrule
1 & 90\% & 80\% & Warm-up; stable GAN convergence \\
2 & 80\% & 70\% & Moderate difficulty; best perceptual quality \\
3 & 70\% & 60\% & Severe degradation robustness \\
\bottomrule
\end{tabular}
\end{table}


\subsection{Experimental Setup and Metrics}

We evaluate the proposed framework using three standard metrics: Peak Signal-to-Noise Ratio (PSNR), Structural Similarity Index (SSIM), and Natural Image Quality Evaluator (NIQE)~\cite{mittal2012making}. While PSNR and SSIM measure pixel-level fidelity, NIQE is a no-reference perceptual metric where lower values indicate more natural and visually realistic outputs. In real-world SR, NIQE is often more representative of human perception than distortion-based metrics~\cite{wang2021realesrgan}.

\subsection{Quantitative Results}

Table~\ref{tab:quantitative_results} summarizes the performance across multiple benchmarks. The results show that our Siamese framework, particularly Phase 2, significantly improves perceptual quality compared to the baseline Real-ESRGAN and our previous lightweight Student model.

\begin{table*}[t]
\centering
\caption{Comparison of Real-ESRGAN variants on benchmark datasets.}
\label{tab:quantitative_results}
\renewcommand{\arraystretch}{1.1}
\resizebox{\textwidth}{!}{%
\begin{tabular}{llccccccccccccc}
\toprule
\multirow{2}{*}{\textbf{Model}} & \multirow{2}{*}{\makecell{\textbf{Training} \\ \textbf{Dataset}}} 
  & \multicolumn{2}{c}{\textbf{Set5}} 
  & \multicolumn{2}{c}{\textbf{Set14}} 
  & \multicolumn{2}{c}{\textbf{RealSR\_V3 (Canon)}} 
  & \multicolumn{2}{c}{\textbf{RealSR\_V3 (Nikon)}} 
  & \multicolumn{2}{c}{\textbf{BSDS100}} 
  & \multicolumn{2}{c}{\textbf{Urban100}}
  & \multirow{2}{*}{\makecell{\textbf{Avg.} \\ \textbf{Time (s)}}} \\
\cmidrule(lr){3-4}\cmidrule(lr){5-6}\cmidrule(lr){7-8}\cmidrule(lr){9-10}\cmidrule(lr){11-12}\cmidrule(lr){13-14}
 & & \textbf{PSNR/SSIM} & \textbf{NIQE} 
   & \textbf{PSNR/SSIM} & \textbf{NIQE} 
   & \textbf{PSNR/SSIM} & \textbf{NIQE} 
   & \textbf{PSNR/SSIM} & \textbf{NIQE} 
   & \textbf{PSNR/SSIM} & \textbf{NIQE} 
   & \textbf{PSNR/SSIM} & \textbf{NIQE} & \\
\midrule

RealSR-JPEG~\cite{ji2020realworld}
  & Bicubic
  & 24.74/0.698                              & 6.11
  & 23.38/0.602                              & 6.01
  & 25.66/0.720                              & 7.57
  & 25.49/0.698                              & 6.68
  & 19.21/0.440                              & 4.37
  & 17.66/0.490                              & 3.97
  & 1.48 \\

BSRGAN~\cite{zhang2021designing}
  & High-order
  & 24.55/0.706                              & 6.71
  & 23.42/0.624                              & 5.72
  & 25.02/0.748                              & \textbf{5.37}
  & 24.41/0.718                              & 5.68
  & 21.96/0.551                              & 4.51
  & 19.48/0.575                              & 4.08
  & 1.42 \\

SwinIR~\cite{liang2021swinir}
  & High-order
  & \textbf{30.72}/\textbf{0.872}            & 7.28
  & \textbf{27.01}/\textbf{0.754}            & 6.29
  & 25.88/0.749                              & 9.08
  & 25.71/0.731                              & 8.61
  & 19.91/0.563                              & 6.81
  & 16.65/0.529                              & 5.31
  & 2.72 \\

Real-ESRGAN~\cite{wang2021realesrgan}
  & Original
  & 26.31/0.801                              & 7.89
  & 25.18/0.691                              & 5.88
  & 25.84/0.780                              & 7.02
  & 25.19/0.755                              & 6.69
  & \textbf{22.34}/0.568                     & 4.46
  & \textbf{20.12}/\textbf{0.639}            & 4.33
  & 1.55 \\

\makecell[l]{\textbf{Real-ESRGAN} \\ \textbf{student (Ours)}}
  & Our Data
  & 25.19/0.767                              & \textbf{6.54}
  & 23.89/0.694                              & 5.82
  & \textbf{26.77}/\textbf{0.805}            & 6.78
  & \textbf{25.93}/\textbf{0.763}            & 6.25
  & 22.31/0.572                              & 4.82
  & 20.00/0.598                              & 4.05
  & 1.45 \\

\makecell[l]{\textbf{Real-ESRGAN} \\ \textbf{siamese (Ours)} \\ \textbf{Phase 2}}
  & Our Data
  & 23.37/0.718                              & 7.98
  & 22.69/0.642                              & 5.91
  & 23.21/0.748                              & 5.71
  & 22.76/0.713                              & 5.69
  & 21.45/\textbf{0.573}                     & \textbf{4.25}
  & 18.57/0.586                              & \textbf{3.82}
  & 1.72 \\

\makecell[l]{\textbf{Real-ESRGAN} \\ \textbf{siamese (Ours)} \\ \textbf{Phase 3}}
  & Our Data
  & 21.98/0.658                              & 8.00
  & 21.44/0.590                              & \textbf{5.58}
  & 22.76/0.733                              & \textbf{5.64}
  & 22.28/0.695                              & \textbf{5.47}
  & 20.70/0.542                              & 4.28
  & 17.80/0.552                              & 3.93
  & \textbf{1.37} \\

\bottomrule
\end{tabular}%
}
\end{table*}

For instance, on the \textbf{RealSR\_V3 (Canon)} dataset, Siamese Phase 2 achieves a NIQE score of \textbf{5.71}, improving over the baseline Real-ESRGAN score of \textbf{6.88}. Similar improvements are observed on \textbf{Urban100} and \textbf{BSDS100}, where our method consistently produces lower NIQE values while maintaining competitive PSNR and SSIM. Although Phase 3 is trained on more severe degradations, Phase 2 provides the best balance between perceptual realism and structural preservation.

We observe that while the Siamese model improves perceptual quality (NIQE), it sometimes yields lower PSNR/SSIM on synthetic benchmarks compared to the original Real-ESRGAN. This reflects the well-known perception-distortion trade-off~\cite{blau2018perception}, where optimizing for visual realism can lead to deviations in pixel-level fidelity. However, in real-world applications, the gain in visual clarity and structural sharpness significantly outweighs the loss in PSNR.

\subsection{Ablation Study}
To verify the effectiveness of our proposed Siamese curriculum learning strategy, we conduct an ablation study comparing different training configurations.
\begin{itemize}
    \item \textbf{Effect of Siamese Framework:} Comparing the Siamese model with our previous student-only model shows that joint training with feature distillation results in more stable texture restoration and fewer artifacts.
    \item \textbf{Effect of Curriculum Learning:} Models trained without the multi-phase curriculum (directly on 60\% quality) often suffer from training instability and generate spurious noise patterns. The progressive schedule (90\% $\to$ 80\% $\to$ 70\% $\to$ 60\%) enables the network to learn robust features gradually.
    \item \textbf{Loss Weight Selection:} The loss weights ($w_{pixel}=2.0, w_{GAN}=0.5$) were chosen to prioritize structural consistency while maintaining adversarial sharpness. Higher GAN weights led to excessive noise, while lower weights resulted in blurry outputs.
\end{itemize}


\subsection{Qualitative Comparison}

The visual results in Fig.~\ref{fig:qualitative_results} further validate the effectiveness of the proposed method. In the \textbf{RealSR\_V3 (Canon\_004)} sample, traditional methods such as ESRGAN and EDSR generate blurry text regions, while the baseline Real-ESRGAN still shows ringing artifacts around edges. In contrast, our Siamese models (especially Phase 2) restore sharper and more readable characters (``5500 4708'') with fewer artifacts.

\begin{figure*}[!t]
    \centering
    \includegraphics[width=0.98\textwidth,height=2.2cm]{images/visual.png}
    \includegraphics[width=0.98\textwidth,height=2.2cm]{images/visual2.png}
    \includegraphics[width=0.98\textwidth,height=2.2cm]{images/visual3.png}
    \caption{Qualitative comparison of different restoration methods on three benchmark datasets: Urban100, BSD100, and RealSR V3. The proposed Siamese Real-ESRGAN framework, particularly Phase 2, produces sharper structures, clearer textures, and more natural visual quality compared to ESRGAN, EDSR, SwinIR, BSRGAN, and the original Real-ESRGAN.}
    \label{fig:qualitative_results}
\end{figure*}

For \textbf{Urban100 (img\_005)}, the building grid patterns present a challenging high-frequency structure. The original Real-ESRGAN shows visible aliasing and jagged edges, whereas our Siamese framework reconstructs cleaner lines and stronger geometric consistency. On \textbf{BSD100 (306006)}, the proposed model better preserves fine coral textures and fish boundaries without excessive smoothing, leading to more natural visual restoration.

Compared to our previous lightweight Student framework, the current weight-sharing Siamese architecture produces fewer hallucinated artifacts and more stable textures. This demonstrates that dynamic supervision from the Teacher pass and progressive curriculum learning significantly improve robustness under severe real-world degradations.


\section{Conclusion}

We proposed an enhanced Real-ESRGAN framework combining high-order degradation with a weight-sharing Siamese training architecture. Using a three-phase curriculum strategy, our model achieves superior perceptual quality and robustness. Experimental results, including a 17.0\% NIQE improvement on RealSR\_V3, confirm that our joint optimization outperforms baseline Real-ESRGAN while maintaining high efficiency. Despite advances, severe degradation remains challenging. Future work will explore automated tuning and generative priors for more extreme scenarios and mobile optimization.

\bibliographystyle{splncs04}
\bibliography{references}

\end{document}
